\section{Two-circle coverage}

\begin{theorem}[982-M01, 982-M02]
For all sufficiently large \(n\), a strictly convex \(n\)-point set with at least
\(0.99n\) points on the union of two circles has at least \(n/10\) pins attaining
\(\lfloor n/2\rfloor\). If the set is partitioned between two nonconcentric circles
with at least \(n/4\) points in each layer, then at least \(n/2000\) pins exceed
\((1/2+10^{-4})n\). All size thresholds are uniform in the circles.
\end{theorem}

\begin{proof}
Assign each sampled point on the circles to exactly one layer. The proof
uses equalities with one target in each of these disjoint layers.

\paragraph{Mixed equalities.}
For different circles with centers \(d,0\) and radii \(R,S\), use complex
coordinates \(p=d+x,q=d+y,r=z\) for a pin and target on the first circle and
a target on the second. Substituting \(\bar x=R^2/x,\bar y=R^2/y,\bar z=S^2/z\)
in the equality of squared distances and clearing denominators gives
\[
 f=R^2z(x-y)^2+yC,\qquad C=(d+x-z)((dx+R^2)z-S^2x).
\]
As a polynomial in \(y\) it is primitive, since its leading coefficient
is \(R^2z\) and \(C(x,0)=-S^2x(d+x)\ne0\). Its discriminant \(C(C-4R^2xz)\)
has odd valuation at \(d+x-z=0\). On that line the other factor of \(C\)
is \(B(x)=dx^2+(d^2+R^2-S^2)x+dR^2\), which vanishes identically only
for identical circles. Thus \(f\) is irreducible. On the complex arc
\(y=x,z=d+x\), generically \(f_z=-xB(x)\ne0\), \(f_y=0\), \(f_{yy}=2R^2(d+x)\ne0\).
The local graph \(z=g(x,y)\) has \(g_x=1,g_y=0,g_{yy}\ne0\) along \(y=x\).
Hence the rational mixed logarithmic derivative of \(g_x/g_y\) has a
second-order pole there. As in the polynomial-graph proof this excludes
an additive separation. All three variables occur. Separate nondegenerate
Cayley parameter changes reduce the circle parameters to real variables,
preserving irreducibility after removing boundary factors, with bounded
degree. Solymosi--Zahl therefore gives \(O(n^{12/7})\) mixed equal-target
pairs \(J_p\) summed over the circle pins in either layer.

\paragraph{From almost-invariance to a group.}
Call surviving pins good for this step if they fail the desired distance
threshold, lie in a circle layer, are not either center, and satisfy
\(J_p\le n/10000\). There are \(o(n)\) discarded pins due to mixed
equalities and at most two centers. If few pins attain the threshold,
the good set has size at least \(0.889n\) in the coverage problem and
\(0.999n\) in the nonconcentric slack problem.
For each layer of size \(m\), reflection about its center--pin axis has
deficit equal to the number of its targets whose reflected partner is
missing. From the formula for paired distances including zero, the sum
of deficits of the two layers is at most
\(2(d_P(p)+1)-(m_1+m_2)+2J_p\). Thus each deficit is at most
\(\eta=0.011n\) for coverage and \(\eta=0.001n\) for slack.

Fix a good pin \(p_0\) and let \(T\) be the rotations of the layer obtained
by composing each good reflection with its reflection. Then
\(|T|\ge |G|/2\), where \(G\) is the good set, because no line contains
more than two vertices. Each member of \(T\) has deficit at most \(2\eta\);
each of \(H=TT^{-1}\) has deficit at most \(4\eta\). Counting ordered pairs
in the layer gives \(|T|\le m^2/(m-2\eta)\) and \(|H|\le m^2/(m-4\eta)\)
when the denominators are positive.

For coverage, assume each layer has at least \(0.07n\) points. The first
upper bound is at most \(m+3\eta\) since \(m\ge6\eta\), giving
\(0.4115n\le m_i\le0.5885n\). The second upper bound is then at most
\(m+5\eta\le0.6435n\) since \(m\ge20\eta\), and \(0.6435n<3|T|/2\).
For slack, the two layer sizes start at \(n/4\), so the same estimates
give \(0.4965n\le m_i\le0.5035n\) and \(|H|\le0.5085n<3|T|/2\).

In either case \(H\) is a finite subgroup. Indeed \(|T\cap hT|>|T|/2\)
for each \(h\in H\), by translating two copies of \(T\) inside the
difference set. For \(h,k\in H\), a common element of the two more-than-half
subsets implies \(hk\in H\). Thus each layer has a group of order
\(h\ge |G|/2\). If an orbit contains \(b_j\) targets, then
\(\sum b_j^2=\sum_{u\in H}|S\cap uS|\ge h(m-4\eta)\), where \(S\) denotes
the layer directions. Some orbit contains at least \(h(1-4\eta/m)\) points,
exceeding \(0.89h\) for coverage or \(0.98h\) for slack. Every good pin's
squared direction about either center belongs to one fixed phase times
the corresponding group by construction of \(T\).

\paragraph{Distinct centers.}
Normalize the first center to zero and the other to \(d>0\). A first-layer
good pin \(p=Rz\), \(|z|=1\), has second squared direction
\(w=z(Rz-d)/(R-dz)\). Both \(z,w\) are fixed phases times roots of unity.
If \(R\ne d\) the graph has a zero and pole at distinct nonzero finite
arguments. This rules out every nonvertical torsion-binomial component
\(z^aw^b=\omega\) with \(b\ne0\). A vertical component is impossible in
the graph, whose numerator and denominator are coprime. The Newton polygon
of its defining Laurent polynomial has area one, so Beukers--Smyth allows
at most 22 such pins. Each layer actually has a positive linear number
of good pins in both parameter regimes. Applying the argument in both
directions forces each circle's radius to equal \(d\). Their nearly full
orbits ensure each open arc of angular length \(\pi/3\) contains a sample
for large \(n\): there are at least \(h/6-1\) orbit positions and fewer
than \(0.11h\) omissions. After scaling the radii to one, the convex hull
of the second layer contains the disk of radius \(\cos(\pi/6)\) at its
center, by support in every direction. The first circle passes through
that center, and its arc centered there of angular length \(\pi/3\)
contains a first-layer point at distance at most \(2\sin(\pi/12)<\cos(\pi/6)\).
This violates strict convex position and completes the slack result.

\paragraph{Concentric coverage.}
For one circle with \(m\) points, at least \(m-2n/3-14\) pins qualify.
For a semicircle sample, count pins with at least \(\lfloor n/2\rfloor\)
targets farther along one direction. Otherwise the center lies inside
the convex hull. Use the torsion-free projection of the circle direction
group from the majority proof. While the residual circle set of size \(u\)
has at least two extreme fibers and \(u\ge2n/3\), delete the smaller of two
extreme fibers. All its pins qualify since they see at least
\(u-\lceil u/4\rceil\ge\lfloor n/2\rfloor\). If the residual size falls below
\(2n/3\), enough qualifying pins have been counted. Otherwise stop at one
fiber, rotate its coordinates to roots of unity, and write \(s=m-u\) for
deleted qualifying pins and \(t=n-u\). When adding \(t\) targets, those
on the same circle in previously deleted fibers cannot merge with the
original distances, by \(bc=a^2\). Each different-radius target merges with
original distances at at most 23 pins by the torsion count in the majority
proof, and each pair of added targets accounts for at most two merges.
Thus at most \(2t(t+22)/(t+1)\le2t+42\) of the residual circle pins fail.
The total \(g\) of qualifying pins satisfies
\(g\ge s+u-2t-42=3m-2n-2s-42\ge3m-2n-2g-42\).
This proves the one-circle estimate. In the coverage setting it handles
coinciding circles or a layer with fewer than \(0.07n\) points, since the
other then has more than \(0.92n\) points.

If the two concentric circles differ in radius and both layers are large,
their rotations have the same group \(H=\mu_h\). Normalize the direction
of \(p_0\) to one. Each of the nearly full orbits contains a good pin
because \(0.89h>0.111n\). Its phase \(\alpha\) therefore satisfies
\(\alpha^2\in H\), so the two core orbits lie in \(\mu_{2h}\). Each core
surrounds the center, and strict convexity allows at most one point per
ray, so the orbits are the two distinct alternating cosets.
A pin outside this grid has all within-core distances distinct on each
core: its reflection does not preserve either orbit. Mixed coincidences
between the two cores at that pin are at most 44, by the same five-term
Laurent equation and Newton polygon area two as in the majority argument.
It consequently sees more than \(1.78h-44>n/2\) distances for sufficiently
large \(n\). Discard the \(k\le0.01n\) off-circle targets and \(\ell<n/10\)
off-grid targets in the two layers for an initial reduced set, of size
\(n-K\) with \(K=k+\ell\le0.11n\).
The two cores have fewer than \(0.11h\) gaps each. There is an inner core
vertex with its two outer neighbors present, since the omissions exclude
fewer than \(0.33h\) positions. This vertex must lie outside the chord
of the neighbors, giving radius ratio \(\rho>\cos(\pi/h)\).

Put \(\tau=\pi/h\) and scale the larger radius to one. From a pin of
radius \(a\in[\rho,1]\), the angular bin \(j\tau\) gives distances
strictly increasing in target radius. Bins \(j,j+1\) are separated for
\(1\le j\le h-2\): at \(\rho=\cos\tau\), the difference between the
minimum squared distance in the later bin and maximum in the earlier
one equals \(2a\sin\tau\sin((j+1)\tau)-\sin^2\tau>0\).
This difference increases with \(\rho\) on the relevant interval.
An outer point without antipode therefore qualifies within the reduced set.
If it has an antipode, that antipode qualifies because the last bin has the
outer point alone, with the strictly greatest distance. Taking each outer
point or its antipode by this rule gives distinct pins, more than \(0.89h\),
each with at least \(\lfloor(n-K)/2\rfloor\) reduced-set distances.

On adding the discarded targets, a new target can merge with the old
two layers at at most 176 selected pins. For unequal target radii, for
each of two pin radii and two target radii the five-term equation in
the pin direction and target relative direction has at most 44 torsion
solutions. An off-grid target in one layer cannot merge with a grid target
in its own layer because reflection at a grid pin preserves the grid.
Pairs of additional targets create at most two equalities each over all
pins. Hence total losses are at most \(K(K+175)\). When \(K>0\), at most
\(2(K+175)\) pins can lose more than \(K/2\). All other selected pins reach
\(\lfloor(n-K)/2\rfloor+K-\lfloor K/2\rfloor\ge\lfloor n/2\rfloor\).
When \(K=0\) they already reach that bound. For large \(n\), the count
remaining exceeds \(0.89(0.4445n)-0.22n-350>n/10\), a contradiction.
\end{proof}

\begin{proposition}[982-M03, 982-M04]
The union of two complete concentric regular polygons of orders at least
three reaches the half bound whenever it is in strictly convex position.
For arbitrary samples on two concentric circles, every pin reaches the
half bound if the radius ratio satisfies no nonzero polynomial of degree
at most two over the field generated by all cosines of sample direction differences.
\end{proposition}
\begin{proof}
The equal-radius case is cocircular. Otherwise denote the complete
polygons by \(A,B\), of orders \(m\ge k\), and scale the radius of \(A\) to
one. If some pin of \(B\) lies off all symmetry axes of \(A\), it sees the
\(m\) distances to \(A\) distinctly. Otherwise no same-ray pair is possible
by strict convexity, and all directions of \(B\) bisect adjacent directions
of \(A\). Thus \(H=m/k\) is an integer. Put \(h=\pi/m\) and \(f(s)=2\sin(s/2)\).
All cross angles are odd multiples of \(h\); \(A\)'s own chord angles
are even multiples. Write \(R\) for the radius of \(B\).

If \(R<1\), strict convexity at any \(B\) vertex between two \(A\) directions
gives \(R>\cos h\). A cross distance at angle \(s\) is below \(f(s)\)
by the cosine rule and less than \(1-\cos h\) below it by the triangle
inequality. For \(s\le\pi-h\),
\[
 f(s)-f(s-h)\ge f(\pi-h)-f(\pi-2h)>1-\cos h.
\]
The final difference after subtracting \(1-\cos h\) equals
\(4\sin^2(h/4)\cos(h/2)\). Thus only the antipodal cross angle, possible
for odd \(m\), may collide with an \(A\)-chord distance.

If \(H=1\), exchange the roles to make \(A\) outer. At an outer pin
the within-layer and cross bins contribute \(\lfloor m/2\rfloor\) and
\(\lceil m/2\rceil\) distances before merging. If the antipodal distance
merges, it equals \(f(\pi-h)=2\cos(h/2)\), because \(1+R>1+\cos h>
2\cos(3h/2)\), the second largest candidate (or a vacuous candidate
when \(m=3\)). The second inequality follows from
\(c^2-\cos(3h/2)=c(1-c)(4c+3)>0\) for \(c=\cos(h/2)\).
So \(R=2c-1\). At an inner pin rescale the inner radius to one.
The cross distances for \(s\le\pi-2h\) are then strictly above \(f(s)\)
and at most \(1/R-1\) higher. The next possible chord distance is separated
by at least \(f(\pi-h)-f(\pi-2h)=2(1-c)(2c+1)>1/R-1\), since
\((2c+1)(2c-1)>1\). The antipodal cross distance exceeds two.
The inner pin therefore sees \(m\) distances.

For \(H\ge2\), suppose first \(m\ge9\). We claim at most four cross
targets at any \(B\) pin have lengths in the old \(A\)-chord set.
The inner-\(B\) case already gives at most one. If \(R>1\), rotate
the \(B\) pin to the positive horizontal ray. The orientation of it
and the \(A\) vertices at angles \(h,3h\) forces \(R<\cos h/\cos(2h)\):
the chord of those vertices has inward half-plane
\(x\cos(2h)+y\sin(2h)<\cos h\). A cross distance is above \(f(s)\)
and at most \(R-1\) higher. For \(s\le\pi-5h\) the next possible chord
is farther than this, since
\[
 f(\pi-4h)-f(\pi-5h)=4\sin(9h/4)\sin(h/4)>
 \cos h/\cos(2h)-1 .
\]
After cancellation the inequality reads
\(\sin(9h/4)\cos(2h)>\sin(3h/2)\cos(h/4)\).
For \(0<h\le\pi/9\), the sine chord bound on \([0,\pi/4]\) gives
\(\sin(9h/4)>2h\), and \(\cos(2h)>3/4\). The right hand side is at most
\(3h/2\). Only two nonantipodal odd bins remain above \(\pi-5h\) in
each parity, and the antipodal distance when present exceeds two.
There are at most two targets in each bin, proving the claim.

At a pin of \(A\) off the axes of \(B\), the \(k\) cross distances
are distinct, so failure would require at least \(\lceil k/2\rceil\)
collisions with its \(A\) chord distances. The axis congruence is
\(H\mid 2j-1\) for \(A\) direction \(2jh\). There are no on-axis
pins for even \(H\); all pins failing would require
\(m\lceil k/2\rceil>4k\), a contradiction. For odd \(H\ge3\) there
are \(k\) on-axis pins, each with \(\lceil k/2\rceil\) distinct
cross distances by reflection. The two orders have the same parity,
so each on-axis pin must suffer at least one merge to fail.
The combined requirement \((m-k)\lceil k/2\rceil+k\ge5k\)
contradicts the four-per-\(B\)-pin bound.

For \(m<9\) the proper-divisor cases are \((m,k)=(6,3),(8,4)\).
For \((6,3)\) failure would make the three cross distances at an \(A\)
pin precisely the \(A\) chords, with squares \(1,3,4\).
Their sum \(3(1+R^2)=8\) would give a right-angle cross square
\(1+R^2=8/3\), impossible. For \((8,4)\) the two layers together
are centrally symmetric and lie in the minimum disk with an
outer contact diameter, so the support-chain lemma applies.

For the arbitrary-sample claim, scale the radii to \(1,R\).
A mixed coincidence at a pin on the first or second layer,
respectively, implies
\[
 2-2\cos u=1+R^2-2R\cos v,\qquad
 R^2(2-2\cos u)=1+R^2-2R\cos v,
\]
with \(u\) the same-layer angle. Each is a nonzero degree-at-most-two
polynomial over the specified field, excluded by hypothesis. Each
separate layer has distance multiplicity at most two, and without
mixed coincidences the two contributions give the half bound.
\end{proof}
